% ================================================================
%  conteudo/material_metodos.tex
% ================================================================

\section{Material e Métodos}

\subsection{Microrganismo e Condições de Cultivo}

Será utilizada a microalga \textit{Chlorella fusca}, obtida de coleção de
cultura institucional. O pré-inóculo será conduzido em erlenmeyers de 500~mL
contendo 200~mL de meio BG-11 modificado sob iluminação artificial contínua
(150~$\mu$mol fótons m$^{-2}$ s$^{-1}$), temperatura de 25~°C e agitação
orbital a 120~rpm, por 7 dias, até atingir fase exponencial de crescimento.

\subsection{Delineamento Experimental}

O experimento será conduzido em dois tratamentos em triplicata:

\begin{itemize}
    \item \textbf{Tratamento A -- Meio Sintético (controle):} meio BG-11
          preparado com reagentes analíticos conforme Rippka et al.\ (1979),
          contendo NaNO\textsubscript{3} (1,5~g/L) como fonte de nitrogênio
          e K\textsubscript{2}HPO\textsubscript{4} (0,04~g/L) como fonte
          de fósforo;

    \item \textbf{Tratamento B -- Meio com Efluente Industrial:} meio BG-11
          com substituição parcial da fonte de nitrogênio e fósforo por
          efluente industrial centrifugado e esterilizado em autoclave,
          diluído na proporção necessária para ajustar as concentrações de
          N e P às mesmas do Tratamento~A.
\end{itemize}

O cultivo será conduzido em fotobiorreatores tubulares de 2~L,
iluminados por LEDs de espectro branco (100~$\mu$mol fótons m$^{-2}$ s$^{-1}$),
com fotoperíodo de 16:8~h (claro:escuro), temperatura de 25~$\pm$~1~°C,
aeração com ar enriquecido a 2\% de CO\textsubscript{2} e vazão de
0,5~vvm, por 15 dias consecutivos.

\subsection{Monitoramento do Crescimento}

O crescimento celular será acompanhado diariamente por leitura de
absorbância a 680~nm em espectrofotômetro (Shimadzu UV-1800), com
conversão para concentração de biomassa seca (g/L) por meio de curva
padrão previamente construída para a espécie. A cada 3 dias serão
retiradas amostras de 30~mL para determinação direta da biomassa seca
por método gravimétrico (centrifugação a 3.500~rpm por 15~min, lavagem
com água destilada, secagem a 105~°C até massa constante).

\subsection{Extração e Quantificação de Lipídeos Totais}

Ao final do cultivo (dia~15), a biomassa será colhida por centrifugação
(3.500~rpm, 15~min), liofilizada e submetida à extração de lipídeos
totais pelo método de Bligh \& Dyer (1959), utilizando mistura de
clorofórmio:metanol:água (1:2:0,8 v/v). A massa lipídica extraída será
determinada gravimetricamente após evaporação do solvente em rota-evaporador
a 40~°C sob pressão reduzida.

\subsection{Caracterização do Perfil de Ácidos Graxos por GC-MS}

Os lipídeos extraídos serão transesterificados com metanol em presença de
HCl 5\% (v/v) a 80~°C por 90~min, produzindo ésteres metílicos de ácidos
graxos (FAMEs). Os FAMEs serão analisados em cromatógrafo gasoso acoplado
a espectrômetro de massas (Shimadzu GCMS-QP2020), utilizando coluna
capilar RTX-WAX (30~m $\times$ 0,25~mm $\times$ 0,25~$\mu$m), hélio como
gás de arraste (1,0~mL/min) e rampa de temperatura de 60~°C (2~min) a
240~°C (5~°C/min). A identificação dos ácidos graxos será feita por
comparação com a biblioteca NIST~17 e por padrões externos certificados
(Sigma-Aldrich FAME mix).

\subsection{Simulação do Processo HEFA}

A partir das frações molares dos ácidos graxos identificados por GC-MS,
será realizada a simulação estequiométrica do processo HEFA em duas etapas:
(i)~HDO, com rendimento mássico em hidrocarbonetos de 80\% e consumo de
H\textsubscript{2} de 0,04~kg/kg de óleo (PEARLSON et al., 2013); e
(ii)~hidrocraqueamento, com distribuição de produtos em nafta~(20\%),
SAF~(50\%) e diesel verde~(30\%). O poder calorífico inferior (LHV) do SAF
será estimado por média ponderada dos LHVs típicos dos hidrocarbonetos
C8--C16 (TURNS, 2012), e a eficiência energética será calculada conforme a
Equação~\ref{eq:eficiencia}:

\begin{equation}
    \eta = \frac{\sum \dot{m}_i \cdot \text{LHV}_i}{\dot{m}_{\text{H}_2}
    \cdot \text{LHV}_{\text{H}_2}}
    \label{eq:eficiencia}
\end{equation}

em que $\dot{m}_i$ é a vazão mássica de cada produto combustível (kg/h),
$\text{LHV}_i$ é o poder calorífico inferior correspondente (MJ/kg) e
$\dot{m}_{\text{H}_2}$ é a vazão de H\textsubscript{2} consumido (kg/h).

\subsection{Análise Orçamentária}

O custo operacional de cada tratamento será estimado com base nos preços
de mercado dos reagentes, do efluente (custo de transporte e pré-tratamento),
energia elétrica (tarifa industrial vigente), mão de obra e consumíveis de
laboratório. Os resultados serão expressos em R\$/kg de biomassa seca e
em R\$/kg de SAF equivalente, permitindo a comparação direta entre os
dois meios de cultivo e a estimativa da redução de custo propiciada pelo
uso do efluente industrial.

\subsection{Análise Estatística}

Os dados serão analisados por meio de ANOVA de um fator seguida de teste
de Tukey (p~$< 0{,}05$) para comparação das médias entre os tratamentos,
utilizando o software R versão 4.3 (R CORE TEAM, 2024).

